\documentclass[11pt]{amsart}
\usepackage[margin=1in]{geometry}
\usepackage{amsmath,amssymb,amsthm,mathtools}
\usepackage{enumitem}
\usepackage{hyperref}
\hypersetup{colorlinks=true,linkcolor=blue,citecolor=blue,urlcolor=blue}
\usepackage[final]{microtype}
\setlength{\emergencystretch}{3em}

\newcommand{\leanrepo}{https://github.com/HQIV/hqiv-lean/blob/main}
\newcommand{\leanfile}[1]{\href{\leanrepo/#1}{\nolinkurl{#1}}}
\newcommand{\leanid}[1]{\nolinkurl{#1}}

\title[Nucleon binding from composite trace]{Nucleon Binding from HQIV Composite-Trace Weights:\\
Lock-In Proton and Neutron Masses in the Lean Library}
\author{Steven Ettinger Jr}
\address{Independent Researcher}
\email{steven@disregardfiat.tech}
\date{Preprint v1 --- \today}

\newtheorem{theorem}{Theorem}[section]
\newtheorem{proposition}[theorem]{Proposition}
\newtheorem{definition}[theorem]{Definition}
\newtheorem{remark}[theorem]{Remark}

\begin{document}

\begin{abstract}
\noindent\textbf{Position in the publication chain.}
This is the seventh paper in the HQIV preprint series and the
fourth Tier-2 closure note after electroweak mass, SM Lagrangian
packaging, and baryogenesis lock-in
\cite{hqiv-mass-spectrum-paper,hqiv-sm-lagrangian-paper,hqiv-baryogenesis-paper}.
Tier~0--1 fix the discrete spine
\cite{ettinger2026hqiv,hqiv-lightcone-paper,hqiv-so8-paper,hqiv-3d-growth-paper,hqiv-oct-action-paper,hqiv-kirchhoff-paper}.
The companion BBN note \cite{hqiv-bbn-paper} consumes the nucleon and
binding outputs proved here.

\noindent\textbf{Result.}
At lock-in shell $\mathrm{referenceM}=4$ the library defines nucleon masses
as constituent energies minus a shared $8\times8$ composite-trace binding
$E_{\mathrm{bind}}$ from explicit generator diagonals and a triadic carrier
state. Proton and neutron differ only through $uud$ versus $udd$ constituent
sums; isospin splitting is carried by hypercharge bookkeeping
($\Delta m_{\mathrm{iso}}=1\,\mathrm{MeV}$ in the witness normalisation).
Machine-checked identities include positivity, shared-surface splitting, and
alignment with meta-harmonic readouts
(\hyperref[lean:derived-nucleon-pack]{derived nucleon mass pack}).
The \hyperref[lean:neutron-scaffold]{neutron stability scaffold} packages
bonded versus free branches; $\beta$-decay lifetimes and Q-values are
explicit open slots.

\noindent\textbf{Single absolute anchor (quarantined).}
Laboratory proton mass $938.272\,\mathrm{MeV}$ appears only as a lapse-readout
anchor condition for comparison, not as an input to
$E_{\mathrm{bind}}$ or constituent sums.

\noindent\textbf{Patch-ontology reader contract.}
Per the patch-theory reader contract below, binding is a discrete network sum
on the null lattice at shell $m$; it is not obtained as a continuum QCD
potential fit.

\noindent\textbf{Status.}
Proved: composite-trace binding spine, derived proton/neutron masses,
constituent isospin splitting, lapse anchor discharge. Scaffold only:
bonded-neutron skew alignment, weak tipping matrix, $880\,\mathrm{s}$ lifetime.
Not claimed: $\beta$-decay rate derivation, nuclear chart beyond scaffold wells.
\end{abstract}
\maketitle

% Shared reader contract: patch QFT + discrete numerics.
% From papers/*.tex:     % Shared reader contract: patch QFT + discrete numerics.
% From papers/*.tex:     % Shared reader contract: patch QFT + discrete numerics.
% From papers/*.tex:     \input{include/patch_theory_messaging.tex}
% From papers/paper/*.tex: \input{../include/patch_theory_messaging.tex}

\paragraph{Patch QFT and discrete numerics (reader contract).}
HQIV (Horizon-Quantized Informational Vacuum) is formulated as a \emph{patch theory}: quantum fields, actions, and physical readouts are attached to discrete patches---shell indices $m\in\mathbb{N}$, finite spacetime index types (e.g.\ $\mathrm{Fin}\,4$), and horizon-limited charts---rather than to a hypothetical fundamental smooth spacetime obtained as a mesh-refinement or ultraviolet limit.
Accordingly, when this text refers to ``QFT,'' it means \emph{patch QFT}: patching rules, finite measurement ledgers, and variational identities on the discrete spine; it is not the claim that the theory is \emph{defined} by recovering ordinary continuum quantum field theory through such a limit.
The numbers that admit the sharpest statements---mass and coupling ladders, curvature ratios, shell thermodynamics, and rational witnesses in \texttt{HQIV\_LEAN}---are objects of \emph{discrete mathematics} on $\mathbb{N}$ (combinatorial growth, cumulative simplex ledgers) together with the ordered field $\mathbb{R}$ as used in the proof assistant for analysis, bounds, and phenomenological display.
Smooth-manifold or continuum field notation, when it appears, is a \emph{translation layer} for comparison with standard physics literature; it does not supersede the discrete definitions.

\paragraph{A continuum is not required, and in places would destroy these results.}
The discrete patch layer is \emph{not} a numerical approximation to a more fundamental continuum theory awaiting future construction; it is the ontology.  Where the literature treats ``the continuum limit'' as the goal, HQIV treats it as a \emph{calculation approximation} for comparison with standard formulas.  The continuum form is an optional convenience for comparison; the proved discrete statement is what carries the physics.  In several load-bearing places---the curvature imprint $\delta_E(m)$ at shell $m$, the harmonic ladder masses $m_\tau\to m_\mu\to m_e$, the lock-in vev $v=\sqrt{\eta_{\rm paper}\,\Omega_k(m_{\rm lockin};m_{\rm lockin})}$, the proton anchor at $m=\mathrm{referenceM}=4$, the baryon-asymmetry $\eta\propto 6^7\sqrt{3}/(m+1)$, and the rational coefficients $\alpha=3/5$, $\gamma=2/5$, $\alpha_{\rm GUT}=1/42$---taking a literal $m\to\infty$ or sub-Planck continuum limit would \emph{erase} the discrete index that is doing the predictive work.  Continuum forms of these objects are therefore not preferred refinements; they are degenerate, information-losing readouts of the same patch data.

\paragraph{An exception: $\mathfrak{so}(8)$ as a de-facto continuum algebra from a discrete construction.}
Principled exception to the discrete-only stance: the closed Lie algebra $\mathfrak{so}(8) \supseteq G_2\cup\{\Delta\}$ generated from the Fano octonion left-multiplication table and the phase-lift generator $\Delta\in(\mathrm{e}_1,\mathrm{e}_7)$.  Although $\mathfrak{so}(8)$ is a continuous (28-dimensional real) Lie algebra, it is built here entirely from \emph{discrete} ingredients (the seven imaginary octonion units, a finite Fano table, and one $\pi/2$ rotation), and its closure to 28 dimensions is a finite linear-algebra certificate (see the \texttt{Hqiv.SO8Closure} / \texttt{Hqiv.SO8ClosureInterface} stack).  The Lie-group structure is therefore a \emph{de-facto continuum algebra obtained from a discrete construction}: continuous in parameter space, but with no continuum spacetime, no continuum measure, and no UV limit attached.  When this manuscript or its companions write $\exp(t\,\Delta)\in\mathrm{SO}(8)$ or treat $\mathfrak{so}(8)$ rotations as smooth, those are statements internal to a finite-dimensional Lie group whose generators are certified by discrete data; they are \emph{not} a continuum-spacetime commitment and do not require a continuum-QFT scaffolding.

% From papers/paper/*.tex: % Shared reader contract: patch QFT + discrete numerics.
% From papers/*.tex:     \input{include/patch_theory_messaging.tex}
% From papers/paper/*.tex: \input{../include/patch_theory_messaging.tex}

\paragraph{Patch QFT and discrete numerics (reader contract).}
HQIV (Horizon-Quantized Informational Vacuum) is formulated as a \emph{patch theory}: quantum fields, actions, and physical readouts are attached to discrete patches---shell indices $m\in\mathbb{N}$, finite spacetime index types (e.g.\ $\mathrm{Fin}\,4$), and horizon-limited charts---rather than to a hypothetical fundamental smooth spacetime obtained as a mesh-refinement or ultraviolet limit.
Accordingly, when this text refers to ``QFT,'' it means \emph{patch QFT}: patching rules, finite measurement ledgers, and variational identities on the discrete spine; it is not the claim that the theory is \emph{defined} by recovering ordinary continuum quantum field theory through such a limit.
The numbers that admit the sharpest statements---mass and coupling ladders, curvature ratios, shell thermodynamics, and rational witnesses in \texttt{HQIV\_LEAN}---are objects of \emph{discrete mathematics} on $\mathbb{N}$ (combinatorial growth, cumulative simplex ledgers) together with the ordered field $\mathbb{R}$ as used in the proof assistant for analysis, bounds, and phenomenological display.
Smooth-manifold or continuum field notation, when it appears, is a \emph{translation layer} for comparison with standard physics literature; it does not supersede the discrete definitions.

\paragraph{A continuum is not required, and in places would destroy these results.}
The discrete patch layer is \emph{not} a numerical approximation to a more fundamental continuum theory awaiting future construction; it is the ontology.  Where the literature treats ``the continuum limit'' as the goal, HQIV treats it as a \emph{calculation approximation} for comparison with standard formulas.  The continuum form is an optional convenience for comparison; the proved discrete statement is what carries the physics.  In several load-bearing places---the curvature imprint $\delta_E(m)$ at shell $m$, the harmonic ladder masses $m_\tau\to m_\mu\to m_e$, the lock-in vev $v=\sqrt{\eta_{\rm paper}\,\Omega_k(m_{\rm lockin};m_{\rm lockin})}$, the proton anchor at $m=\mathrm{referenceM}=4$, the baryon-asymmetry $\eta\propto 6^7\sqrt{3}/(m+1)$, and the rational coefficients $\alpha=3/5$, $\gamma=2/5$, $\alpha_{\rm GUT}=1/42$---taking a literal $m\to\infty$ or sub-Planck continuum limit would \emph{erase} the discrete index that is doing the predictive work.  Continuum forms of these objects are therefore not preferred refinements; they are degenerate, information-losing readouts of the same patch data.

\paragraph{An exception: $\mathfrak{so}(8)$ as a de-facto continuum algebra from a discrete construction.}
Principled exception to the discrete-only stance: the closed Lie algebra $\mathfrak{so}(8) \supseteq G_2\cup\{\Delta\}$ generated from the Fano octonion left-multiplication table and the phase-lift generator $\Delta\in(\mathrm{e}_1,\mathrm{e}_7)$.  Although $\mathfrak{so}(8)$ is a continuous (28-dimensional real) Lie algebra, it is built here entirely from \emph{discrete} ingredients (the seven imaginary octonion units, a finite Fano table, and one $\pi/2$ rotation), and its closure to 28 dimensions is a finite linear-algebra certificate (see the \texttt{Hqiv.SO8Closure} / \texttt{Hqiv.SO8ClosureInterface} stack).  The Lie-group structure is therefore a \emph{de-facto continuum algebra obtained from a discrete construction}: continuous in parameter space, but with no continuum spacetime, no continuum measure, and no UV limit attached.  When this manuscript or its companions write $\exp(t\,\Delta)\in\mathrm{SO}(8)$ or treat $\mathfrak{so}(8)$ rotations as smooth, those are statements internal to a finite-dimensional Lie group whose generators are certified by discrete data; they are \emph{not} a continuum-spacetime commitment and do not require a continuum-QFT scaffolding.


\paragraph{Patch QFT and discrete numerics (reader contract).}
HQIV (Horizon-Quantized Informational Vacuum) is formulated as a \emph{patch theory}: quantum fields, actions, and physical readouts are attached to discrete patches---shell indices $m\in\mathbb{N}$, finite spacetime index types (e.g.\ $\mathrm{Fin}\,4$), and horizon-limited charts---rather than to a hypothetical fundamental smooth spacetime obtained as a mesh-refinement or ultraviolet limit.
Accordingly, when this text refers to ``QFT,'' it means \emph{patch QFT}: patching rules, finite measurement ledgers, and variational identities on the discrete spine; it is not the claim that the theory is \emph{defined} by recovering ordinary continuum quantum field theory through such a limit.
The numbers that admit the sharpest statements---mass and coupling ladders, curvature ratios, shell thermodynamics, and rational witnesses in \texttt{HQIV\_LEAN}---are objects of \emph{discrete mathematics} on $\mathbb{N}$ (combinatorial growth, cumulative simplex ledgers) together with the ordered field $\mathbb{R}$ as used in the proof assistant for analysis, bounds, and phenomenological display.
Smooth-manifold or continuum field notation, when it appears, is a \emph{translation layer} for comparison with standard physics literature; it does not supersede the discrete definitions.

\paragraph{A continuum is not required, and in places would destroy these results.}
The discrete patch layer is \emph{not} a numerical approximation to a more fundamental continuum theory awaiting future construction; it is the ontology.  Where the literature treats ``the continuum limit'' as the goal, HQIV treats it as a \emph{calculation approximation} for comparison with standard formulas.  The continuum form is an optional convenience for comparison; the proved discrete statement is what carries the physics.  In several load-bearing places---the curvature imprint $\delta_E(m)$ at shell $m$, the harmonic ladder masses $m_\tau\to m_\mu\to m_e$, the lock-in vev $v=\sqrt{\eta_{\rm paper}\,\Omega_k(m_{\rm lockin};m_{\rm lockin})}$, the proton anchor at $m=\mathrm{referenceM}=4$, the baryon-asymmetry $\eta\propto 6^7\sqrt{3}/(m+1)$, and the rational coefficients $\alpha=3/5$, $\gamma=2/5$, $\alpha_{\rm GUT}=1/42$---taking a literal $m\to\infty$ or sub-Planck continuum limit would \emph{erase} the discrete index that is doing the predictive work.  Continuum forms of these objects are therefore not preferred refinements; they are degenerate, information-losing readouts of the same patch data.

\paragraph{An exception: $\mathfrak{so}(8)$ as a de-facto continuum algebra from a discrete construction.}
Principled exception to the discrete-only stance: the closed Lie algebra $\mathfrak{so}(8) \supseteq G_2\cup\{\Delta\}$ generated from the Fano octonion left-multiplication table and the phase-lift generator $\Delta\in(\mathrm{e}_1,\mathrm{e}_7)$.  Although $\mathfrak{so}(8)$ is a continuous (28-dimensional real) Lie algebra, it is built here entirely from \emph{discrete} ingredients (the seven imaginary octonion units, a finite Fano table, and one $\pi/2$ rotation), and its closure to 28 dimensions is a finite linear-algebra certificate (see the \texttt{Hqiv.SO8Closure} / \texttt{Hqiv.SO8ClosureInterface} stack).  The Lie-group structure is therefore a \emph{de-facto continuum algebra obtained from a discrete construction}: continuous in parameter space, but with no continuum spacetime, no continuum measure, and no UV limit attached.  When this manuscript or its companions write $\exp(t\,\Delta)\in\mathrm{SO}(8)$ or treat $\mathfrak{so}(8)$ rotations as smooth, those are statements internal to a finite-dimensional Lie group whose generators are certified by discrete data; they are \emph{not} a continuum-spacetime commitment and do not require a continuum-QFT scaffolding.

% From papers/paper/*.tex: % Shared reader contract: patch QFT + discrete numerics.
% From papers/*.tex:     % Shared reader contract: patch QFT + discrete numerics.
% From papers/*.tex:     \input{include/patch_theory_messaging.tex}
% From papers/paper/*.tex: \input{../include/patch_theory_messaging.tex}

\paragraph{Patch QFT and discrete numerics (reader contract).}
HQIV (Horizon-Quantized Informational Vacuum) is formulated as a \emph{patch theory}: quantum fields, actions, and physical readouts are attached to discrete patches---shell indices $m\in\mathbb{N}$, finite spacetime index types (e.g.\ $\mathrm{Fin}\,4$), and horizon-limited charts---rather than to a hypothetical fundamental smooth spacetime obtained as a mesh-refinement or ultraviolet limit.
Accordingly, when this text refers to ``QFT,'' it means \emph{patch QFT}: patching rules, finite measurement ledgers, and variational identities on the discrete spine; it is not the claim that the theory is \emph{defined} by recovering ordinary continuum quantum field theory through such a limit.
The numbers that admit the sharpest statements---mass and coupling ladders, curvature ratios, shell thermodynamics, and rational witnesses in \texttt{HQIV\_LEAN}---are objects of \emph{discrete mathematics} on $\mathbb{N}$ (combinatorial growth, cumulative simplex ledgers) together with the ordered field $\mathbb{R}$ as used in the proof assistant for analysis, bounds, and phenomenological display.
Smooth-manifold or continuum field notation, when it appears, is a \emph{translation layer} for comparison with standard physics literature; it does not supersede the discrete definitions.

\paragraph{A continuum is not required, and in places would destroy these results.}
The discrete patch layer is \emph{not} a numerical approximation to a more fundamental continuum theory awaiting future construction; it is the ontology.  Where the literature treats ``the continuum limit'' as the goal, HQIV treats it as a \emph{calculation approximation} for comparison with standard formulas.  The continuum form is an optional convenience for comparison; the proved discrete statement is what carries the physics.  In several load-bearing places---the curvature imprint $\delta_E(m)$ at shell $m$, the harmonic ladder masses $m_\tau\to m_\mu\to m_e$, the lock-in vev $v=\sqrt{\eta_{\rm paper}\,\Omega_k(m_{\rm lockin};m_{\rm lockin})}$, the proton anchor at $m=\mathrm{referenceM}=4$, the baryon-asymmetry $\eta\propto 6^7\sqrt{3}/(m+1)$, and the rational coefficients $\alpha=3/5$, $\gamma=2/5$, $\alpha_{\rm GUT}=1/42$---taking a literal $m\to\infty$ or sub-Planck continuum limit would \emph{erase} the discrete index that is doing the predictive work.  Continuum forms of these objects are therefore not preferred refinements; they are degenerate, information-losing readouts of the same patch data.

\paragraph{An exception: $\mathfrak{so}(8)$ as a de-facto continuum algebra from a discrete construction.}
Principled exception to the discrete-only stance: the closed Lie algebra $\mathfrak{so}(8) \supseteq G_2\cup\{\Delta\}$ generated from the Fano octonion left-multiplication table and the phase-lift generator $\Delta\in(\mathrm{e}_1,\mathrm{e}_7)$.  Although $\mathfrak{so}(8)$ is a continuous (28-dimensional real) Lie algebra, it is built here entirely from \emph{discrete} ingredients (the seven imaginary octonion units, a finite Fano table, and one $\pi/2$ rotation), and its closure to 28 dimensions is a finite linear-algebra certificate (see the \texttt{Hqiv.SO8Closure} / \texttt{Hqiv.SO8ClosureInterface} stack).  The Lie-group structure is therefore a \emph{de-facto continuum algebra obtained from a discrete construction}: continuous in parameter space, but with no continuum spacetime, no continuum measure, and no UV limit attached.  When this manuscript or its companions write $\exp(t\,\Delta)\in\mathrm{SO}(8)$ or treat $\mathfrak{so}(8)$ rotations as smooth, those are statements internal to a finite-dimensional Lie group whose generators are certified by discrete data; they are \emph{not} a continuum-spacetime commitment and do not require a continuum-QFT scaffolding.

% From papers/paper/*.tex: % Shared reader contract: patch QFT + discrete numerics.
% From papers/*.tex:     \input{include/patch_theory_messaging.tex}
% From papers/paper/*.tex: \input{../include/patch_theory_messaging.tex}

\paragraph{Patch QFT and discrete numerics (reader contract).}
HQIV (Horizon-Quantized Informational Vacuum) is formulated as a \emph{patch theory}: quantum fields, actions, and physical readouts are attached to discrete patches---shell indices $m\in\mathbb{N}$, finite spacetime index types (e.g.\ $\mathrm{Fin}\,4$), and horizon-limited charts---rather than to a hypothetical fundamental smooth spacetime obtained as a mesh-refinement or ultraviolet limit.
Accordingly, when this text refers to ``QFT,'' it means \emph{patch QFT}: patching rules, finite measurement ledgers, and variational identities on the discrete spine; it is not the claim that the theory is \emph{defined} by recovering ordinary continuum quantum field theory through such a limit.
The numbers that admit the sharpest statements---mass and coupling ladders, curvature ratios, shell thermodynamics, and rational witnesses in \texttt{HQIV\_LEAN}---are objects of \emph{discrete mathematics} on $\mathbb{N}$ (combinatorial growth, cumulative simplex ledgers) together with the ordered field $\mathbb{R}$ as used in the proof assistant for analysis, bounds, and phenomenological display.
Smooth-manifold or continuum field notation, when it appears, is a \emph{translation layer} for comparison with standard physics literature; it does not supersede the discrete definitions.

\paragraph{A continuum is not required, and in places would destroy these results.}
The discrete patch layer is \emph{not} a numerical approximation to a more fundamental continuum theory awaiting future construction; it is the ontology.  Where the literature treats ``the continuum limit'' as the goal, HQIV treats it as a \emph{calculation approximation} for comparison with standard formulas.  The continuum form is an optional convenience for comparison; the proved discrete statement is what carries the physics.  In several load-bearing places---the curvature imprint $\delta_E(m)$ at shell $m$, the harmonic ladder masses $m_\tau\to m_\mu\to m_e$, the lock-in vev $v=\sqrt{\eta_{\rm paper}\,\Omega_k(m_{\rm lockin};m_{\rm lockin})}$, the proton anchor at $m=\mathrm{referenceM}=4$, the baryon-asymmetry $\eta\propto 6^7\sqrt{3}/(m+1)$, and the rational coefficients $\alpha=3/5$, $\gamma=2/5$, $\alpha_{\rm GUT}=1/42$---taking a literal $m\to\infty$ or sub-Planck continuum limit would \emph{erase} the discrete index that is doing the predictive work.  Continuum forms of these objects are therefore not preferred refinements; they are degenerate, information-losing readouts of the same patch data.

\paragraph{An exception: $\mathfrak{so}(8)$ as a de-facto continuum algebra from a discrete construction.}
Principled exception to the discrete-only stance: the closed Lie algebra $\mathfrak{so}(8) \supseteq G_2\cup\{\Delta\}$ generated from the Fano octonion left-multiplication table and the phase-lift generator $\Delta\in(\mathrm{e}_1,\mathrm{e}_7)$.  Although $\mathfrak{so}(8)$ is a continuous (28-dimensional real) Lie algebra, it is built here entirely from \emph{discrete} ingredients (the seven imaginary octonion units, a finite Fano table, and one $\pi/2$ rotation), and its closure to 28 dimensions is a finite linear-algebra certificate (see the \texttt{Hqiv.SO8Closure} / \texttt{Hqiv.SO8ClosureInterface} stack).  The Lie-group structure is therefore a \emph{de-facto continuum algebra obtained from a discrete construction}: continuous in parameter space, but with no continuum spacetime, no continuum measure, and no UV limit attached.  When this manuscript or its companions write $\exp(t\,\Delta)\in\mathrm{SO}(8)$ or treat $\mathfrak{so}(8)$ rotations as smooth, those are statements internal to a finite-dimensional Lie group whose generators are certified by discrete data; they are \emph{not} a continuum-spacetime commitment and do not require a continuum-QFT scaffolding.


\paragraph{Patch QFT and discrete numerics (reader contract).}
HQIV (Horizon-Quantized Informational Vacuum) is formulated as a \emph{patch theory}: quantum fields, actions, and physical readouts are attached to discrete patches---shell indices $m\in\mathbb{N}$, finite spacetime index types (e.g.\ $\mathrm{Fin}\,4$), and horizon-limited charts---rather than to a hypothetical fundamental smooth spacetime obtained as a mesh-refinement or ultraviolet limit.
Accordingly, when this text refers to ``QFT,'' it means \emph{patch QFT}: patching rules, finite measurement ledgers, and variational identities on the discrete spine; it is not the claim that the theory is \emph{defined} by recovering ordinary continuum quantum field theory through such a limit.
The numbers that admit the sharpest statements---mass and coupling ladders, curvature ratios, shell thermodynamics, and rational witnesses in \texttt{HQIV\_LEAN}---are objects of \emph{discrete mathematics} on $\mathbb{N}$ (combinatorial growth, cumulative simplex ledgers) together with the ordered field $\mathbb{R}$ as used in the proof assistant for analysis, bounds, and phenomenological display.
Smooth-manifold or continuum field notation, when it appears, is a \emph{translation layer} for comparison with standard physics literature; it does not supersede the discrete definitions.

\paragraph{A continuum is not required, and in places would destroy these results.}
The discrete patch layer is \emph{not} a numerical approximation to a more fundamental continuum theory awaiting future construction; it is the ontology.  Where the literature treats ``the continuum limit'' as the goal, HQIV treats it as a \emph{calculation approximation} for comparison with standard formulas.  The continuum form is an optional convenience for comparison; the proved discrete statement is what carries the physics.  In several load-bearing places---the curvature imprint $\delta_E(m)$ at shell $m$, the harmonic ladder masses $m_\tau\to m_\mu\to m_e$, the lock-in vev $v=\sqrt{\eta_{\rm paper}\,\Omega_k(m_{\rm lockin};m_{\rm lockin})}$, the proton anchor at $m=\mathrm{referenceM}=4$, the baryon-asymmetry $\eta\propto 6^7\sqrt{3}/(m+1)$, and the rational coefficients $\alpha=3/5$, $\gamma=2/5$, $\alpha_{\rm GUT}=1/42$---taking a literal $m\to\infty$ or sub-Planck continuum limit would \emph{erase} the discrete index that is doing the predictive work.  Continuum forms of these objects are therefore not preferred refinements; they are degenerate, information-losing readouts of the same patch data.

\paragraph{An exception: $\mathfrak{so}(8)$ as a de-facto continuum algebra from a discrete construction.}
Principled exception to the discrete-only stance: the closed Lie algebra $\mathfrak{so}(8) \supseteq G_2\cup\{\Delta\}$ generated from the Fano octonion left-multiplication table and the phase-lift generator $\Delta\in(\mathrm{e}_1,\mathrm{e}_7)$.  Although $\mathfrak{so}(8)$ is a continuous (28-dimensional real) Lie algebra, it is built here entirely from \emph{discrete} ingredients (the seven imaginary octonion units, a finite Fano table, and one $\pi/2$ rotation), and its closure to 28 dimensions is a finite linear-algebra certificate (see the \texttt{Hqiv.SO8Closure} / \texttt{Hqiv.SO8ClosureInterface} stack).  The Lie-group structure is therefore a \emph{de-facto continuum algebra obtained from a discrete construction}: continuous in parameter space, but with no continuum spacetime, no continuum measure, and no UV limit attached.  When this manuscript or its companions write $\exp(t\,\Delta)\in\mathrm{SO}(8)$ or treat $\mathfrak{so}(8)$ rotations as smooth, those are statements internal to a finite-dimensional Lie group whose generators are certified by discrete data; they are \emph{not} a continuum-spacetime commitment and do not require a continuum-QFT scaffolding.


\paragraph{Patch QFT and discrete numerics (reader contract).}
HQIV (Horizon-Quantized Informational Vacuum) is formulated as a \emph{patch theory}: quantum fields, actions, and physical readouts are attached to discrete patches---shell indices $m\in\mathbb{N}$, finite spacetime index types (e.g.\ $\mathrm{Fin}\,4$), and horizon-limited charts---rather than to a hypothetical fundamental smooth spacetime obtained as a mesh-refinement or ultraviolet limit.
Accordingly, when this text refers to ``QFT,'' it means \emph{patch QFT}: patching rules, finite measurement ledgers, and variational identities on the discrete spine; it is not the claim that the theory is \emph{defined} by recovering ordinary continuum quantum field theory through such a limit.
The numbers that admit the sharpest statements---mass and coupling ladders, curvature ratios, shell thermodynamics, and rational witnesses in \texttt{HQIV\_LEAN}---are objects of \emph{discrete mathematics} on $\mathbb{N}$ (combinatorial growth, cumulative simplex ledgers) together with the ordered field $\mathbb{R}$ as used in the proof assistant for analysis, bounds, and phenomenological display.
Smooth-manifold or continuum field notation, when it appears, is a \emph{translation layer} for comparison with standard physics literature; it does not supersede the discrete definitions.

\paragraph{A continuum is not required, and in places would destroy these results.}
The discrete patch layer is \emph{not} a numerical approximation to a more fundamental continuum theory awaiting future construction; it is the ontology.  Where the literature treats ``the continuum limit'' as the goal, HQIV treats it as a \emph{calculation approximation} for comparison with standard formulas.  The continuum form is an optional convenience for comparison; the proved discrete statement is what carries the physics.  In several load-bearing places---the curvature imprint $\delta_E(m)$ at shell $m$, the harmonic ladder masses $m_\tau\to m_\mu\to m_e$, the lock-in vev $v=\sqrt{\eta_{\rm paper}\,\Omega_k(m_{\rm lockin};m_{\rm lockin})}$, the proton anchor at $m=\mathrm{referenceM}=4$, the baryon-asymmetry $\eta\propto 6^7\sqrt{3}/(m+1)$, and the rational coefficients $\alpha=3/5$, $\gamma=2/5$, $\alpha_{\rm GUT}=1/42$---taking a literal $m\to\infty$ or sub-Planck continuum limit would \emph{erase} the discrete index that is doing the predictive work.  Continuum forms of these objects are therefore not preferred refinements; they are degenerate, information-losing readouts of the same patch data.

\paragraph{An exception: $\mathfrak{so}(8)$ as a de-facto continuum algebra from a discrete construction.}
Principled exception to the discrete-only stance: the closed Lie algebra $\mathfrak{so}(8) \supseteq G_2\cup\{\Delta\}$ generated from the Fano octonion left-multiplication table and the phase-lift generator $\Delta\in(\mathrm{e}_1,\mathrm{e}_7)$.  Although $\mathfrak{so}(8)$ is a continuous (28-dimensional real) Lie algebra, it is built here entirely from \emph{discrete} ingredients (the seven imaginary octonion units, a finite Fano table, and one $\pi/2$ rotation), and its closure to 28 dimensions is a finite linear-algebra certificate (see the \texttt{Hqiv.SO8Closure} / \texttt{Hqiv.SO8ClosureInterface} stack).  The Lie-group structure is therefore a \emph{de-facto continuum algebra obtained from a discrete construction}: continuous in parameter space, but with no continuum spacetime, no continuum measure, and no UV limit attached.  When this manuscript or its companions write $\exp(t\,\Delta)\in\mathrm{SO}(8)$ or treat $\mathfrak{so}(8)$ rotations as smooth, those are statements internal to a finite-dimensional Lie group whose generators are certified by discrete data; they are \emph{not} a continuum-spacetime commitment and do not require a continuum-QFT scaffolding.


\section{Scope and claim discipline}
\label{sec:scope}

The mass-spectrum note \cite{hqiv-mass-spectrum-paper} establishes the
electroweak ladder and lock-in discipline; the SM Lagrangian bridge
\cite{hqiv-sm-lagrangian-paper} packages sector projections; baryogenesis
\cite{hqiv-baryogenesis-paper} fixes $\eta$ normalisation without deriving
nucleon spectroscopy. \textbf{This note answers: how are proton and neutron
masses defined from the same $8\times8$ network binding, and what is proved
versus scaffold for free-neutron decay?}

\begin{enumerate}[leftmargin=2em]
\item \textbf{Proved in Lean (binding spine).}
  $E_{\mathrm{bind}}=\sum_k w_k\cdot(\texttt{latticeSimplexCount}\cdot\alpha_{\mathrm{eff}})$
  from composite traces (\hyperref[lean:e-bind-network]{network binding formula});
  shared binding at $\mathrm{referenceM}$ from the nucleon trace witness
  (\hyperref[lean:nucleon-shared-binding]{nucleon shared binding});
  $\texttt{derivedProtonMass}$ and $\texttt{derivedNeutronMass}$ as constituent
  minus shared binding (\hyperref[lean:derived-nucleon-pack]{derived nucleon
  mass pack}); $\texttt{derivedDeltaM}$ equals constituent isospin splitting
  (\hyperref[lean:isospin-splitting]{constituent isospin splitting theorem}).
\item \textbf{Proved in Lean (readout alignment).}
  Raw hadron masses from the same composite trace recover
  \texttt{derivedProtonMass} / \texttt{derivedNeutronMass}
  (\hyperref[lean:composite-trace-mass]{composite-trace mass identities});
  proton lapse anchor condition discharges laboratory comparison when assumed
  (\hyperref[lean:proton-anchor]{proton anchor condition}).
\item \textbf{Scaffold (neutron stability).}
  Bonded versus free neutron embeddings, overlap energy
  $\Delta E_{\mathrm{overlap}}$, and $\Gamma=\Delta E/\hbar$ width identity
  (\hyperref[lean:neutron-scaffold]{neutron stability scaffold}); Conjecture
  $\beta$ slots (skew alignment, weak matrix) \emph{not} discharged.
\item \textbf{Downstream consumer.}
  BBN network weights and cluster $Q$ values import the same composite trace
  \cite{hqiv-bbn-paper}; this note does not develop abundances.
\item \textbf{Not claimed.}
  Derivation of $880\,\mathrm{s}$ free-neutron lifetime; full nuclear chart;
  MS\=bar current-quark fits as inputs.
\end{enumerate}

\section{The $8\times 8$ composite-trace binding law}
\label{sec:binding}

The \hyperref[lean:bound-states]{bound-states module} makes the network
picture explicit. Generator diagonals $\mathrm{diag}_k(i)$ and an eight-component
state $\psi_i$ define a trace weight per $\mathfrak{so}(8)$ generator:
\[
  w_k = \sum_{i=0}^{7} \mathrm{diag}_k(i)\,\psi_i^2,
  \qquad
  E_{\mathrm{bind}}(m) = \sum_k w_k \cdot \underbrace{(\texttt{latticeSimplexCount}(m)\cdot\alpha_{\mathrm{eff}}(m))}_{\text{shell coupling}}.
\]
The \hyperref[lean:e-bind-composite]{composite-trace binding definition}
packages $w$ from $(\mathrm{diag},\psi)$ so every hadron model instantiates the
same spine rather than a separate potential ansatz.

For nucleons the \hyperref[lean:nucleon-trace-witness]{nucleon trace witness}
uses a triadic state on three carrier slots of one generator family. The
theorem \hyperref[lean:triple-coupling]{triadic triple coupling at shell $m$}
shows this witness equals three times the shell coupling at $m$---a concrete
$8\times8$ normalisation used throughout BBN and excited-state readouts.

\section{Derived proton and neutron masses at lock-in}
\label{sec:masses}

Fix $\mathrm{referenceM}=4$ as the proton lock-in shell
\cite{hqiv-mass-spectrum-paper,hqiv-kirchhoff-paper}. The
\hyperref[lean:derived-nucleon-mass]{derived nucleon mass module} defines:
\begin{align*}
  E_{\mathrm{shared}} &= \texttt{nucleonSharedBinding\_MeV}
    = E_{\mathrm{bind}}(\mathrm{referenceM};\mathrm{diag}_{\mathrm{nuc}},\psi_{\mathrm{nuc}}),\\
  M_p^{\mathrm{derived}} &= E_{\mathrm{constituent}}(uud) - E_{\mathrm{shared}},\\
  M_n^{\mathrm{derived}} &= E_{\mathrm{constituent}}(udd) - E_{\mathrm{shared}},\\
  \Delta M &= M_n^{\mathrm{derived}} - M_p^{\mathrm{derived}}
    = E_{\mathrm{constituent}}(udd)-E_{\mathrm{constituent}}(uud).
\end{align*}
Thus the neutron--proton mass difference is \emph{not} an independent fit knob:
it is forced by constituent bookkeeping once the shared binding is shared.

The hypercharge sign flip in the meta-resonance layer sets
$\Delta m_{\mathrm{iso}}=1\,\mathrm{MeV}$ in the witness normalisation
(\hyperref[lean:isospin-gap-one]{isospin gap unit theorem}). This is distinct
from $\Delta M$ above: the scaffold overlap energy for a free neutron uses
$\Delta m_{\mathrm{iso}}$ as a bookkeeping gap, not the full $\beta$ Q-value
($\sim 0.78\,\mathrm{MeV}$).

Positivity $M_p^{\mathrm{derived}}, M_n^{\mathrm{derived}}>0$ and alignment with
meta-harmonic readouts are proved in the same module
(\hyperref[lean:derived-nucleon-pack]{derived nucleon mass pack}).

\section{Lapse readout and the proton anchor}
\label{sec:lapse}

The \hyperref[lean:lapse-mass-readout]{lapse mass readout module} expresses
lab comparison through shell lapse $\Phi$:
\[
  M_{\mathrm{lab}} = \frac{M_{\mathrm{raw}}}{\texttt{shellLapse}(\mathrm{referenceM},\Phi,t)}.
\]
The \hyperref[lean:proton-anchor]{proton anchor condition} states that when
$\texttt{shellLapse}\cdot 938.272 = M_p^{\mathrm{derived}}$, the lapse readout
returns $938.272\,\mathrm{MeV}$. This is a \textbf{one-point calibration witness}
for comparison to CODATA/PDG tables, not a second simultaneous fit of binding
and constituents.

\section{Bonded versus free neutrons (scaffold)}
\label{sec:neutron}

The \hyperref[lean:neutron-scaffold]{neutron stability scaffold} contrasts:

\begin{itemize}[leftmargin=2em]
\item \textbf{Bonded branch.}
  Nuclear embedding with well depth, $\Omega_k$-type readout $\ge 1$, and
  skew-alignment slot (Conjecture $\beta$). When the curvature ledger is closed,
  overlap energy reduces to the isospin gap
  (\hyperref[lean:overlap-closed-ledger]{overlap equals isospin gap}).
\item \textbf{Free branch.}
  Sub-lock-in readout $\omega<1$ adds a curvature deficit to overlap energy;
  strong-width scaffold $\Gamma=\Delta E/\hbar$ is proved as an identity, not
  matched to $880\,\mathrm{s}$.
\end{itemize}

Continuous-$\xi$ lock-in calibration
(\hyperref[lean:xi-lockin]{$\xi$-lock-in calibration}) and bonded readout at
$\mathrm{referenceM}$ are proved packaging lemmas linking the scaffold to
baryogenesis normalisation \cite{hqiv-baryogenesis-paper}.

\section{What this note does not claim}
\label{sec:not-claimed}

\begin{itemize}[leftmargin=2em]
\item Weak-process matrix elements or $\beta$-decay half-lives.
\item Identification of $\Delta E_{\mathrm{overlap}}$ with the full $\beta$
  Q-value.
\item Nuclear binding energies beyond the effective-well slots in
  \hyperref[lean:nuclear-spectra]{nuclear spectra module}.
\item Any continuum-limit or lattice-spacing extrapolation programme.
\end{itemize}

\section{Conclusion}

HQIV nucleon masses at lock-in are \emph{constituent sums minus a shared
composite-trace binding} on the $\mathfrak{so}(8)$ network, with isospin
splitting carried by constituent hypercharge bookkeeping. The proof graph is
sorry-free through positivity, splitting identities, and lapse anchor discharge;
free-neutron $\beta$ decay remains an explicit scaffold for a later upgrade.
The BBN companion \cite{hqiv-bbn-paper} imports the same binding spine for
light-element $Q$ values and partition weights.

\section*{Acknowledgments}

Machine-checked proofs live in the public HQIV Lean repository \cite{hqiv-lean}.

\appendix

\section{Lean module index}
\label{app:lean-index}

\begin{description}[leftmargin=2.5em,style=nextline]
\item[\hypertarget{lean:bound-states}{Bound states / network spine}]
  \leanfile{Hqiv/Physics/BoundStates.lean} ---
  \leanid{E\_bind\_from\_network}, \leanid{E\_bind\_from\_composite\_trace}.
\item[\hypertarget{lean:e-bind-network}{Network binding formula}]
  \leanid{E\_bind\_from\_network} in \leanfile{Hqiv/Physics/BoundStates.lean}.
\item[\hypertarget{lean:e-bind-composite}{Composite-trace binding definition}]
  \leanid{E\_bind\_from\_composite\_trace} in \leanfile{Hqiv/Physics/BoundStates.lean}.
\item[\hypertarget{lean:nucleon-trace-witness}{Nucleon trace witness}]
  \leanfile{Hqiv/Physics/QuarkMetaResonance.lean} ---
  \leanid{nucleonTraceDiagonal}, \leanid{nucleonTraceState}.
\item[\hypertarget{lean:nucleon-shared-binding}{Nucleon shared binding}]
  \leanid{nucleonSharedBinding\_MeV}, \leanid{nucleonSharedBinding\_from\_composite\_trace}.
\item[\hypertarget{lean:triple-coupling}{Triadic triple coupling}]
  \leanid{E\_bind\_nucleon\_trace\_eq\_triple\_coupling} in
  \leanfile{Hqiv/Physics/MetaHorizonExcitedStates.lean}.
\item[\hypertarget{lean:derived-nucleon-mass}{Derived nucleon mass module}]
  \leanfile{Hqiv/Physics/DerivedNucleonMass.lean}.
\item[\hypertarget{lean:derived-nucleon-pack}{Derived nucleon mass pack}]
  \leanid{derivedProtonMass\_pos}, \leanid{derivedNeutronMass\_pos},
  \leanid{derivedProtonMass\_eq\_meta\_harmonics},
  \leanid{proton\_mass\_from\_shared\_harmonics}.
\item[\hypertarget{lean:isospin-splitting}{Constituent isospin splitting}]
  \leanid{constituent\_isospin\_splitting} in
  \leanfile{Hqiv/Physics/DerivedNucleonMass.lean}.
\item[\hypertarget{lean:isospin-gap-one}{Isospin gap unit theorem}]
  \leanid{nucleonIsospinGap\_eq\_one} in
  \leanfile{Hqiv/Physics/QuarkMetaResonance.lean}.
\item[\hypertarget{lean:composite-trace-mass}{Composite-trace mass identities}]
  \leanid{proton\_raw\_hadron\_mass\_from\_composite\_trace},
  \leanid{nucleonSharedBinding\_uses\_composite\_trace\_only} in
  \leanfile{Hqiv/Physics/LapseMassReadout.lean}.
\item[\hypertarget{lean:lapse-mass-readout}{Lapse mass readout module}]
  \leanfile{Hqiv/Physics/LapseMassReadout.lean}.
\item[\hypertarget{lean:proton-anchor}{Proton anchor condition}]
  \leanid{ProtonAnchorCondition}, \leanid{proton\_anchor\_condition\_discharge}.
\item[\hypertarget{lean:neutron-scaffold}{Neutron stability scaffold}]
  \leanfile{Hqiv/Physics/NeutronBindingStabilityScaffold.lean}.
\item[\hypertarget{lean:overlap-closed-ledger}{Overlap equals isospin gap}]
  \leanid{freeNeutronOverlapEnergy\_eq\_isospinGap\_when\_ledger\_closed}.
\item[\hypertarget{lean:xi-lockin}{$\xi$-lock-in calibration}]
  \leanid{bondedNeutronReadoutCalibrated} in
  \leanfile{Hqiv/Physics/NeutronBindingStabilityScaffold.lean}.
\item[\hypertarget{lean:nuclear-spectra}{Nuclear spectra module}]
  \leanfile{Hqiv/Physics/NuclearAndAtomicSpectra.lean}.
\item[\hypertarget{lean:hqiv-nuclei}{Isotope valley scaffold}]
  \leanfile{Hqiv/Physics/HQIVNuclei.lean}.
\end{description}

\bibliographystyle{plain}
\bibliography{../references}

\end{document}
